% !PRINTED 622
\weberpage{622}
%<<CONT
thus carries the five pairs over into five others; and since there
are only ten pairs of five things, all pairs occur in the two
sequences (3).

One can now form half-metacyclic functions in very different
ways. The simplest is
\weq{4}{u = x_0x_1 + x_1x_2 + x_2x_3 + x_3x_4 + x_4x_0,}
which by $t$ passes into
\weq{5}{u' = x_0x_2 + x_2x_4 + x_4x_1 + x_1x_3 + x_3x_0}
and $u'$ itself belongs to the half-metacyclic functions.

If
\weq{6}{f(x) = x^5 - ax^4 + bx^3 - cx^2 + dx - e = 0}
is the equation whose roots are $x_0, x_1, x_2, x_3, x_4$, then accordingly
\weq{7}{u + u' = b.}

The function
\weq{8}{y = u - u'}
is likewise half-metacyclic, but the square $y^2$ is full metacyclic
and therefore the root of a resolvent of degree 6,
as \gesperrt{Lagrange} already found.

If $\sqrt{\Delta}$ is the product of the ten root differences $(x_0 - x_1)$,
$(x_0 - x_2) \dots$, so that $\Delta$ is the discriminant of $f(x)$, then
\weq{9}{Y = \dfrac{u - u'}{\sqrt{\Delta}}}

is likewise full metacyclic and the root of an equation of degree 6 $^{1)}$.

In order to form the conjugate functions to $u$, we denote
by $M$ the full metacyclic, by $N$ the half-metacyclic,
by $A$ the alternating group, and decompose $A$ into the subsidiary groups:
\weq{10}{\begin{aligned}A &= N + N(1,2)(3,4) + Nt(0,1) + Nt(0,2)\\ &+ Nt(0,3) + Nt(0,4).\end{aligned}}

These subsidiary groups are in fact all different from one another.
For if, e.\,g., $N = N(1,2)(3,4)$, then $(1,2)(3,4)$ would have to
occur in $N$, and hence
\[ (1,2)(3,4)t = (1,2)(3,4)(1,2,4,3) = (1,4) \]
would have to occur in $M$, and this is not possible, because in $M$ no
permutation appears that leaves two figures unchanged; and in

\par\noindent\rule{4cm}{0.4pt}\par
\noindent $^{1)}$ \gesperrt{Jacobi}, ``observatiunculae ad theoriam aequationum pertinentes'',
Crelle's Journ. Bd. 13. \gesperrt{Jacobi's Werke}, Bd. 3. \gesperrt{Cayley}, philos. transactions
 1861, Collected math. papers, Vol. IV, p. 309.
%CONT>>
